\chapter{Học để làm gì khi đời không dễ}
\markboth{Học để làm gì khi đời không dễ}{Học để làm gì khi đời không dễ}

\section{Ra đời còn chưa chắc kiếm được việc, học dữ vậy để làm gì?}
\begin{truyen}{Ra đời còn chưa chắc kiếm được việc, học dữ vậy để làm gì?}{Classroom Talk}
\chuhoa{H}{ọc sinh: ``Nghe nói tốt nghiệp xong vẫn thất nghiệp đầy. Vậy học cắm đầu giờ để làm gì hả thầy?''}
\textit{Student: ``People say many graduates still end up unemployed. So why should we bury ourselves in studying now?''}\\[4pt]

\teacherpair{Không ai dám hứa học xong là đời dễ. Nhưng không học thì lúc đời khó, em còn ít đồ nghề hơn để tự cứu mình.}{No one can promise life will become easy after school. But if you do not learn, you will have fewer tools to help yourself when life gets hard.}

\studentpair{Vậy học không phải vé thắng, mà là đỡ tay trắng hả thầy?}{So studying is not a winning ticket. It just keeps me from facing life empty-handed, right?}

\teacherpair{Đúng. Học không cứu em khỏi mọi khổ. Nhưng nó làm em bớt bất lực.}{Exactly. Studying does not save you from every hardship. But it makes you less helpless.}
\end{truyen}

\begin{baihoc}
Không phải cứ học là chắc thắng. Nhưng không chịu học thì thua thường đến sớm hơn.
\textit{(Study does not guarantee success, but refusing to learn often brings failure sooner.)}
\end{baihoc}

\begin{ghinhoanh}
\textbf{Short English to remember:}

Study is not a magic ticket.

It is a tool against helplessness.
\end{ghinhoanh}

\begin{tuhocnhanh}
1. Em đang học để thắng, hay học để đỡ bất lực? \textit{(Do you study to win, or to be less helpless?)}

2. Khi gặp chuyện khó, em có đủ kỹ năng để tự xoay chưa? \textit{(When life gets hard, do you have enough skills to handle it?)}
\end{tuhocnhanh}
\ngancach

\section{Nhà em nói kiếm tiền mới quan trọng, học sau cũng được}
\begin{truyen}{Nhà em nói kiếm tiền mới quan trọng, học sau cũng được}{Classroom Talk}
\chuhoa{H}{ọc sinh: ``Nhà em lo tiền từng ngày. Nên nghe ai nói học cho cao siêu em thấy xa lắm. Kiếm tiền mới là việc gấp.''}
\textit{Student: ``My family worries about money every day. So when people talk about lofty learning goals, it feels far away. Earning money feels more urgent.''}\\[4pt]

\teacherpair{Vậy em muốn cả đời chỉ đủ sống từng ngày, hay sau này đỡ bị đồng tiền dắt đi?}{So do you want to survive day by day forever, or do you want to be less controlled by money later on?}

\studentpair{Chắc ai nghèo cũng nghĩ tiền gấp trước đã.}{I guess anyone who grows up poor would think urgent money comes first.}

\teacherpair{Đúng. Nhưng nếu chỉ lo cái gấp mà bỏ cái gốc, vài năm nữa em vẫn chạy vòng đó thôi.}{True. But if you only chase what is urgent and neglect the foundation, you will still be running in circles a few years from now.}
\end{truyen}

\begin{baihoc}
Người khó khăn càng cần học thực chất, vì đó là cách chậm mà bền để không nghèo mãi.
\textit{(Those in hardship need real learning even more, because it is the slower but steadier way out.)}
\end{baihoc}

\begin{ghinhoanh}
\textbf{Short English to remember:}

Urgent money matters.

But skills decide how long money stays.
\end{ghinhoanh}

\begin{tuhocnhanh}
1. Em đang giải quyết cái gấp hay đang xây cái bền? \textit{(Are you solving the urgent, or building the lasting?)}

2. Em muốn học gì để sau này kiếm tiền đỡ bấp bênh hơn? \textit{(What do you want to learn so your future income is less unstable?)}
\end{tuhocnhanh}
\ngancach

\section{Em học hoài vẫn đuối, vậy cố để làm gì?}
\begin{truyen}{Em học hoài vẫn đuối, vậy cố để làm gì?}{Classroom Talk}
\chuhoa{H}{ọc sinh: ``Bạn khác học một hiểu mười. Em học năm hiểu một. Cứ kiểu này thì cố để làm gì?''}
\textit{Student: ``Other people study once and understand ten things. I study five times and understand one. If it stays like this, why should I keep trying?''}\\[4pt]

\teacherpair{Em muốn bằng người khác ngay, hay muốn em của tháng sau đỡ tệ hơn em của tháng này?}{Do you want to be equal to other people right away, or do you want the version of you next month to be less weak than the version of you now?}

\studentpair{Nghe vậy thì đỡ cay hơn, nhưng vẫn mệt.}{That sounds less bitter, but it still feels exhausting.}

\teacherpair{Mệt là thật. Nhưng mệt vì đang đi lên khác với mệt vì đứng yên rồi tự chê mình.}{The exhaustion is real. But being tired because you are moving upward is different from being tired because you are standing still and insulting yourself.}
\end{truyen}

\begin{baihoc}
Cố gắng không phải để chứng minh mình hơn ai. Nhiều khi chỉ là để mai đỡ dở hơn hôm nay.
\textit{(Effort is not always about beating others. Sometimes it is just about being less weak tomorrow.)}
\end{baihoc}

\begin{ghinhoanh}
\textbf{Short English to remember:}

Slow progress is still progress.

Do not compare speed before checking direction.
\end{ghinhoanh}

\begin{tuhocnhanh}
1. Em đang thua người khác hay đang bỏ cuộc với chính mình? \textit{(Are you losing to others, or quitting on yourself?)}

2. Một điều nhỏ nào chứng minh em đã khá hơn trước? \textit{(What small thing proves you are better than before?)}
\end{tuhocnhanh}
\ngancach

\section{Em học chỉ để ba mẹ đỡ chửi, vậy có sao không?}
\begin{truyen}{Em học chỉ để ba mẹ đỡ chửi, vậy có sao không?}{Classroom Talk}
\chuhoa{H}{ọc sinh: ``Nói thật nha thầy, nhiều bữa em học không phải vì thích. Em học để khỏi bị la, khỏi bị so với con nhà người ta.''}
\textit{Student: ``Honestly, many times I do not study because I like it. I study so I do not get scolded or compared with someone else’s child.''}\\[4pt]

\teacherpair{Lý do đó không đẹp, nhưng ít ra còn kéo em ngồi vào bàn. Vấn đề là em có chịu đi tiếp xa hơn lý do đó không.}{That reason is not beautiful, but at least it gets you to sit down and study. The question is whether you are willing to move beyond that reason.}

\studentpair{Nghĩa là sao?}{What do you mean by that?}

\teacherpair{Ban đầu em học vì sợ. Sau đó em phải dần học vì tương lai của chính em. Nếu đứng yên ở nỗi sợ, em sẽ học trong bực bội suốt.}{At first, fear may be what gets you started. But later, you need to learn for your own future. If you stay stuck in fear, you will study in resentment forever.}
\end{truyen}

\begin{baihoc}
Động lực xấu vẫn tốt hơn không có động lực. Nhưng nếu không nâng nó lên, em sẽ rất nhanh kiệt sức.
\textit{(A weak motivation is still better than none, but if it never grows, burnout comes quickly.)}
\end{baihoc}

\begin{ghinhoanh}
\textbf{Short English to remember:}

Fear can start a habit.

Purpose is what keeps it alive.
\end{ghinhoanh}

\begin{tuhocnhanh}
1. Lý do học hiện tại của em là gì? \textit{(What is your current reason for studying?)}

2. Em muốn đổi từ học vì sợ sang học vì điều gì? \textit{(What do you want to study for besides fear?)}
\end{tuhocnhanh}
\ngancach